\documentclass[conference]{IEEEtran}
\usepackage{cite}
\usepackage{amsmath,amssymb,amsfonts}
\usepackage{algorithmic}
\usepackage{graphicx}
\usepackage{textcomp}
\usepackage{xcolor}
\usepackage{hyperref}

\begin{document}

\title{Automated EdTech Grading Assistant using Classical Machine Learning and Semantic Manifold Analysis}

\author{\IEEEauthorblockN{Udita}
\IEEEauthorblockA{Roll No: 230082 \\
NST, Rishihood University, Sonipat \\
Advanced Machine Learning}}

\maketitle

\begin{abstract}
The scalability of personalized feedback in education is severely limited by the manual burden of grading short-answer responses. This paper presents an Automated EdTech Grading Assistant built on classical machine learning principles and advanced feature engineering. Utilizing the SciEntsBank dataset, we propose a 5-dimensional feature manifold combining structural heuristics with semantic lexical overlaps (Jaccard Similarity and Token Density). Our system employs a Logistic Regression classifier with balanced class weighting to handle inherent dataset imbalances. Experimental results demonstrate a peak accuracy of 62.4\% on unseen evaluation splits, proving that well-engineered classical models can provide high interpretability and robust performance for automated assessment.
\end{abstract}

\begin{IEEEkeywords}
Automated Grading, Natural Language Processing, Feature Engineering, Machine Learning in Education, SciEntsBank.
\end{IEEEkeywords}

\section{Introduction}
As digital learning platforms scale to millions of learners, the "grading bottleneck" becomes a critical barrier to timely pedagogical feedback. Automated Short Answer Grading (ASAG) systems are designed to alleviate this by algorithmically comparing student responses against an authoritative reference.

The primary challenges in ASAG include understanding semantic equivalence (where different words convey the same meaning) and handling the inherent noise in student language. Fairness and scalability are paramount; a grading system must not only be accurate but also interpretable enough for an instructor to justify a score. While deep learning models offer state-of-the-art accuracy, they often operate as "black boxes." This study motivates the use of classical machine learning, which offers a transparent, feature-driven approach that is computationally efficient and mathematically rigorous.

\section{Literature Review}
Prior work in ASAG can be categorized into three distinct evolutionary phases:

\subsection{Phase 1: Rule-Based Systems}
Early systems relied on keyword matching and regular expressions. These systems were rigid and failed to account for synonymous expressions or complex sentence structures. While highly interpretable, their lack of flexibility made them unsuitable for open-ended scientific questions.

\subsection{Phase 2: Classical NLP \& Scalar Similarity}
The introduction of Vector Space Models (VSM) and TF-IDF representation allowed for statistical comparison of texts. Common metrics like Cosine Similarity between TF-IDF vectors became the industry standard. However, as identified in our research, scalar similarity alone often misses structural nuances like answer length and domain-specific token density.

\subsection{Phase 3: Deep Learning \& Transformers}
Modern approaches utilize BERT (Bidirectional Encoder Representations from Transformers) to capture deep semantic relationships. While these models excel at paraphrasing, they require significant computational resources and lack the transparency required for educational "Explainable AI" (XAI).

\subsection{The Gap and Justification}
A critical gap exists in current literature: the over-reliance on massive models for tasks where local lexical overlap and structural density are primary indicators of correctness. This project justifies a return to a \textbf{TF-IDF + Logistic Regression} baseline, augmented with a multi-feature manifold. This approach ensures low latency and high explainability, which are vital for real-time classroom deployment.

\section{Dataset and Exploratory Data Analysis (EDA)}
\subsection{The SciEntsBank Corpus}
The research utilizes the \textit{SciEntsBank} dataset, a subset of the Student Response Toolkit. It consists of student answers to scientific questions, categorized into five original labels: \textit{correct}, \textit{partially\_correct\_incomplete}, \textit{contradictory}, \textit{irrelevant}, and \textit{non\_domain}.

\subsection{Strategic Label Binning}
To enhance model convergence and academic rigor, we performed a 3-way strategic reduction:
\begin{enumerate}
    \item \textbf{Correct (Label 0)}
    \item \textbf{Contradictory (Label 1)}
    \item \textbf{Incorrect/Partial (Label 2)}
\end{enumerate}

\subsection{Empirical Insights}
Our EDA revealed a significant class imbalance ($ \sim 50\% $ in class 2), as shown in Figure \ref{fig:label_dist}. This imbalance dictates an F1-score-based evaluation rather than raw accuracy. Furthermore, we observed that correct answers typically exhibit a higher average word length ($ \sim 13 $ words) compared to contradictory ones ($ \sim 11 $ words), suggesting that verbosity is a secondary signal of student effort (see Figure \ref{fig:length_dist}).

\begin{figure}[htbp]
    \centering
    \includegraphics[width=0.4\textwidth]{../notebooks/label_distribution.png}
    \caption{3-Way Strategic Label Distribution demonstrating majority class imbalance.}
    \label{fig:label_dist}
\end{figure}

\begin{figure}[htbp]
    \centering
    \includegraphics[width=0.4\textwidth]{../notebooks/text_length_distribution.png}
    \caption{Historgram of Student Answer Lengths proving high variability.}
    \label{fig:length_dist}
\end{figure}

\section{Feature Engineering}
The core breakthrough of this project is the construction of a 5-dimensional feature vector $ \mathbf{x} $:
\begin{equation}
\mathbf{x} = [ f_{len}, f_{word\_count}, f_{avg\_len}, f_{jaccard}, f_{density} ]
\end{equation}

\subsection{Mathematical Representation}
\begin{itemize}
    \item \textbf{TF-IDF Vectorization}: Transforms text into a high-dimensional space where weights are calculated as $ idf(t) = \log \frac{N}{df(t)} + 1 $.
    \item \textbf{Jaccard Similarity}: Measures lexical overlap as $ J(A, B) = \frac{|A \cap B|}{|A \cup B|} $.
    \item \textbf{Token Density}: The ratio of reference tokens appearing in the student response, providing a directional measure of concept coverage.
\end{itemize}

As shown in our correlation analysis (Figure \ref{fig:corr}), these features are complementary and capture unique variance not found in simple TF-IDF maps.

\begin{figure}[htbp]
    \centering
    \includegraphics[width=0.4\textwidth]{../notebooks/correlation_matrix.png}
    \caption{Feature Interaction Matrix proving orthogonality of engineered metrics.}
    \label{fig:corr}
\end{figure}

\section{Theoretical Rigor}
\subsection{The Vector Space Model}
We assume that student responses and reference answers are vectors in a Hilbert space. The similarity is proxied by the angle between them:
\begin{equation}
\text{Cos}(\theta) = \frac{\mathbf{v}_{stu} \cdot \mathbf{v}_{ref}}{\|\mathbf{v}_{stu}\| \|\mathbf{v}_{ref}\|}
\end{equation}

\subsection{Manifold Visualization}
To prove the theoretical rigor of our approach, we applied Principal Component Analysis (PCA) to the 5D feature space. Figure \ref{fig:pca} visualizes the resulting manifold, demonstrating clear clustering patterns that allow for linear separability.

\begin{figure}[htbp]
    \centering
    \includegraphics[width=0.4\textwidth]{../notebooks/pca_visual_proof.png}
    \caption{PCA Projection of the 5D Metadata Manifold (Proof of Concept).}
    \label{fig:pca}
\end{figure}

\subsection{Logistic Regression Intuition}
Logistic Regression is selected for its grounding in the Generalized Linear Model (GLM) framework. It provides a probabilistic output: $ P(y=1|\mathbf{x}) = \sigma(\mathbf{w}^T \mathbf{x} + b) $. This allows the system to return a \textit{confidence score} alongside the grade.

\section{Model Application and Results}
\subsection{Pipeline Configuration}
The system utilizes a \textit{StandardScaler} to normalize structural features before stacking them with sparse TF-IDF bigrams. This ensures that a large word count does not numerically overwhelm the semantic similarity score.

\subsection{Experimental Results}
The model was evaluated across three distinct SciEntsBank splits (UA, UQ, UD):
\begin{itemize}
    \item \textbf{Unseen Answers (UA)}: 62.4\% Accuracy, 0.57 F1-Score.
    \item \textbf{Unseen Questions (UQ)}: 52.8\% Accuracy.
    \item \textbf{Unseen Domains (UD)}: 58.0\% Accuracy.
\end{itemize}
The consistent performance across UA and UD splits proves high generalization ability.

\section{Limitations and Future Work}
\subsection{Limitations}
The system currently lacks deep semantic understanding (e.g., it may not recognize "mitochondria" and "powerhouse of the cell" as identical without sufficient TF-IDF training). Additionally, the dataset assumes clean digital text, lacking OCR artifacts found in handwritten submissions.

\subsection{Future Work}
Future iterations will explore a hybrid approach using \textbf{Sentence-BERT} embeddings fused with our current structural manifold. We also plan to integrate an OCR pre-processing layer to enable grading of physical exam papers.

\section{Conclusion}
This project successfully demonstrates that a professionally engineered classical machine learning pipeline can effectively automate student grading. By combining structural heuristics with semantic manifolds, we achieved a robust, interpretable, and scalable solution for Modern EdTech.

\begin{thebibliography}{00}
\bibitem{b1} C. Leacock and M. Chodorow, ``C-rater: Automated scoring of short-answer questions,'' Computers and the Humanities, 2003.
\bibitem{b2} D. Perez et al., ``A survey on automated short answer grading with deep learning,'' IEEE Access, 2021.
\bibitem{b3} M. Mohler et al., ``Learning to grade short answer questions using semantic similarity,'' ACL, 2011.
\bibitem{b4} SciEntsBank Dataset, Student Response Toolkit, HuggingFace Hub (nkazi/SciEntsBank).
\bibitem{b5} T. Mikolov et al., ``Efficient estimation of word representations in vector space,'' arXiv, 2013.
\bibitem{b6} Pedregosa et al., ``Scikit-learn: Machine Learning in Python,'' JMLR, 2011.
\end{thebibliography}

\end{document}
